\answer{ex:18:1}{(a)~$0.9949\le w\le1.0049$, 즉 $|1-w|\lesssim0.005$. (b)~시냅스 $K\ge400$개.}
\answer{ex:18:2}{$(\omega_R\tau_w)^2=\sqrt{\alpha(\alpha+2+2\varrho)}/\varrho-1$, $\varrho=\tau_m/\tau_w$; $f_R\approx11.1$~Hz; $|Z(\omega_R)|/|Z(0)|=4.6$.}
\answer{ex:18:3}{(a)~$f=1/\bigl(\tau\ln\frac{\mu}{\mu-\theta}\bigr)$; $1$~Hz에는 $\mu/\theta-1\approx3.7\cdot10^{-44}$이 필요하다. (b)~$T=\pi\tau/\sqrt\mu$, $f\approx3.2$~Hz: $f\propto\sqrt\mu$.}
\answer{ex:18:4}{적분기는 $f\lesssim17.7$~Hz에서 발화하고(저역 통과 필터), 공진기는 $5.5$--$22$~Hz 대역에서만 발화한다(대역 통과 필터).}
\answer{ex:18:5}{(a)~$T_{\min}=\frac{\tau}{w-1}\ln\frac{0.5}{0.3}\approx10.2\ms$. (b)~$B>w-1-0.2=0.3$.}
\answer{ex:18:6}{조정의 시상수는 $\tau/(\eta\langle r^2\rangle)$, $\eta=\tau\ln10/60\,\text{s}\approx3.8\cdot10^{-3}$; $I\neq0$이면 가중치가 치우치므로 $e$에서 명령의 사본 $I/\tau$를 빼야 한다.}
\answer{ex:18:7}{열린 문제. 공진기는 $11$~Hz에서 $1.35$배만 유리하고 $1$~Hz에서는 $17$배 불리하다. 재구성의 주된 이득은 문턱값에 의한 대역 선택에 있다(연습문제~\ref{ex:18:4}).}
